\scanpage{29}
\mspage{(10)}
\margin{\uline{NB} Ce sous-groupoïde $\mathrm{III}_{J,i}$ (de groupe fondamental $\widehat{\mathbb Z}$) engendré par les sommets $Q_j$ $(j\in J\setminus\{i\})$ et les deux $\unsure{\widetilde a_i^\omega|\text{arêtes}}$ \unsure{orientés} qui les joignent, \uuline{n'est \deleted{\unsure{donc}} pas} stable \uline{par $u$} !}
en d'autres termes faisons opérer $u$ sur $\pi(Q_{\omega'i})$, identifié à $\widehat\Pi_{0,3}$ via $\psi_{\omega'i}^{\omega'}$, $\widetilde\lambda_i^\omega$ identifié à $l_0$, on trouve que $u(l_0)$ est conjugué de $l_0^p$, \deleted{\unsure{avec}} la conjugaison étant effectuée par $\eta^{-1}$ (par la formule $(30)$).

Le passage de $u$ à $\tau(u)$ s'exprime par
\deleted{\unsure{J'en prends} maintenant les points $P_\omega^{\unsure{\omega}}$ $(\omega\in\uline{\omega})$}
le passage de $\widetilde\alpha$ à
\[
(31)\qquad\widetilde\alpha'=\varepsilon_0\tau_\infty(\widetilde\alpha^{-1})
={\underbrace{\rho^{-1}\underline{\widetilde\sigma_\infty\vphantom{(\widetilde\alpha^{-1})}}}_{\widetilde\sigma_0}}\underline{(\widetilde\alpha^{-1})\vphantom{\widetilde\sigma_\infty}}
\]
\margin{où $\widetilde\sigma_\infty=\sigma_\infty\tau_\infty=\tau_\infty\sigma_\infty$ est dans l'image de l'hom canonique
\diagram{\begin{tikzcd}[ampersand replacement=\&]
\unsure{\widetilde G_3|G_3}\arrow[r]\arrow[d,phantom,"\simeq"{rotate=90,description}] \& \widetilde\Pi_{0,3}\\
\deleted{\text{\unsure{du}}\ \unsure{\widetilde G_3|G_3}} \&
\end{tikzcd}}
stabilisateur de $P_+$ dans $G_3\times\{1,\tau\}$}
\[
\widetilde\sigma_i(l_i)=l_i^{-1}
\]
\[
\widetilde\sigma_i(l_j)=l_k^{-1}\quad\text{si }\{0,1,\infty\}=\{i,j,k\}
\]
\uline{NB} Le groupe $\widetilde G_3$ opère sur $\Pi_{0,3}$ et contient \deleted{$\unsure{\xi_3|\widetilde G_3}$} \deleted{\unsure{dont il}} $\{1,\rho^{-1},\rho^{-2}_{\unsure{\prime\prime|\mathrm{II}}}\}$ comme ss-groupe \inserted{$\widetilde G_3^+$} d'indice 2, \deleted{\unsure{et|dont} $\widetilde G_3$ \unsure{conjugue} $\rho$} $\widetilde G_3/\widetilde G_3^+$ opère sur $\widetilde G_3^+$ en échangeant $\rho$ et $\rho^{-1}$, il s'ensuit que $\widetilde G_3$ opère sur $V$ ens des solutions $x$ de l'équation
\[
x\rho^{-1}(x)\rho^{-2}(x)=1
\]
par
\[
(32)\qquad\Theta_g x=g.x^{\operatorname{sg}(g)}
\]
et on a ici
\[
(33)\qquad\widetilde\alpha'=\Theta_{\widetilde\sigma_0}(\widetilde\alpha)
\]
\margin{28}
